% Revised against the implemented OpenNWM-finalLAM-100k and navigation-LAM
% checkpoints on 2026-09-23.  The two figures in the supplied draft are
% intentionally omitted because their image contents were not available for
% verification; do not restore their placeholder captions.

\section{Experiments}
\label{sec:experiments}

\subsection{Experimental Setup}

\paragraph{Datasets.}
The world-model batches in OpenNWM's two NWM training stages draw from
disjoint source sets.  Stage~1 uses the 13-source \textsc{NavAnywhere}
action-free video corpus.  Stage~2 uses the training splits of RECON, HuRoN
(stored under the historical \texttt{sacson} key in our implementation),
SCAND, and TartanDrive.  Go Stanford is excluded from both stages and is used
only to evaluate transfer to an unseen environment.  We retain signed temporal
offsets and do not discard an example merely because its target precedes its
observation.  The standalone LAM benchmark described below uses a different
model and data protocol; we state it separately to avoid conflating LAM
reconstruction with OpenNWM training.

\paragraph{Main-comparison scope and provenance.}
Table~\ref{tab:main_prediction_navigation} contains three distinct comparisons.
Autoregressive visual prediction is evaluated on RECON, direct visual
prediction is evaluated on SACSoN, and navigation is evaluated on RECON.
Visual results are reported at 4 and 16\,s.  These dataset--horizon pairs
define the intended comparison layout, not an already frozen local OpenNWM
protocol.  We compare methods only within an evaluation block: in particular,
the direct and autoregressive rows use different datasets and do not support a
cross-block ranking.  Baseline values in this table are taken from their cited
publications and retain the original training and evaluation protocols; they
are reference points rather than a fully unified re-evaluation.  The published
NWM navigation result, for example, uses a CEM population of 120, whereas our
local planning protocol uses 80 candidates.  Unless stated otherwise, the
OpenNWM row uses the final checkpoint defined in the implementation paragraph
below.

Our canonical local direct-prediction benchmark currently contains 500 fixed
RECON windows at 4\,s and is used for the controlled ablations below.  It does
not provide the SACSoN 4- and 16-s measurements required by the direct block
of Table~\ref{tab:main_prediction_navigation}; those cells remain unreported
until a separate SACSoN protocol is frozen and evaluated.  Autoregressive
rollout uses 150 fixed RECON windows, four context frames, evaluation at 1 and
4\,FPS and at horizons of 1, 2, 4, 8, and 16\,s, and prediction seed 0.
OpenNWM uses 250-step DDPM sampling.  Planning uses 100 fixed cases per
dataset, CEM with a population of 80 and five elites, three stochastic
repetitions, one optimization update, and an eight-step horizon at 0.25\,s per
step.  The planning seed is 42 and the cost is LPIPS-Alex.  Consequently, our
reported values are fixed-seed, single-checkpoint point estimates, not means
and standard deviations over five generated rollouts or multiple independently
trained models.

\paragraph{Controlled NWM-ablation protocol.}
The controlled NWM ablations intentionally use a different training recipe
from the main OpenNWM checkpoint.  TimePT, GeoPT, and LatentPT use the same
CDiT-B/2 backbone and downstream splits and, unless explicitly stated, the
earlier schedule of a 10k-step real-action-adapter warm-up followed by 100k
steps of joint fine-tuning.  Direct prediction is measured on RECON at 4\,s,
rollout is measured on RECON at 1 and 4\,FPS, and OOD direct prediction is
measured on Go Stanford at 4\,s.  The full internal planning benchmark uses
the CEM80 protocol above on RECON and SCAND.  These ablations support
within-table comparisons only; their absolute values should not be compared
to the main-model row, which uses the final 60k-stage-1/3k-warm-up/100k-joint
recipe.

\paragraph{Metrics.}
We use LPIPS-Alex~\citep{zhang2018unreasonable},
DreamSim~\citep{fu2023dreamsim}, and PSNR for locally run visual-prediction
experiments.  The main table additionally retains published FID
values~\citep{heusel2017gans} where supplied by prior work; our current local
benchmark does not compute FID, and we make no FVD claim.  In our local
OpenNWM evaluation, ATE and RPE~\citep{sturm2012benchmark} are translation
RMSEs in waypoint-normalized units, computed without trajectory alignment or
scale correction; RPE uses adjacent poses ($\Delta=1$).  They are evaluated on
the waypoint trajectories selected through model-predicted observations, not
on outputs of a separately trained trajectory-regression head.  The published
navigation baselines retain the metric definitions and protocols of their
source papers.  The LAM reconstruction benchmark uses PSNR and LPIPS as its
primary metrics.

\paragraph{Baselines.}
The autoregressive block compares with NWM and AR Forcing, the direct block
with NWM and RAE-NWM, and the navigation block with GNM, NoMaD, NWM, and LAM.
DIAMOND is not part of the reported main table.  The standalone LAM benchmark
separately evaluates the released DreamDojo LAM~\citep{gao2026dreamdojo},
villa-X~\citep{chen2025villa0x0}, DiLA~\citep{zhang2026dila},
Olaf-World~\citep{jiang2026olaf}, and the MoTo tokenizer released with
LARA~\citep{liu2026lara,chen2024moto}.  We use ``public DreamDojo LAM'' for
that baseline to distinguish it from our navigation LAM implementation in the
DreamDojo codebase.

\paragraph{OpenNWM implementation.}
We instantiate OpenNWM with a CDiT-B/2 diffusion transformer (approximately
197M parameters), rather than CDiT-XL, and encode $224\times224$ images with
the frozen \texttt{stabilityai/sd-vae-ft-ema} tokenizer.  Each observation has
four context frames and four sampled target frames.  Stage~1 conditions the
world model on layer-normalized, 32-dimensional proxies precomputed by our
60k-step Pixel+Action navigation LAM.  We sample signed frame offsets in
$[-64,64]$ and condition every sample on relative time.  When
$|\Delta t|\leq 8$, the valid pair additionally receives its exact LAM proxy;
larger-offset samples have no proxy but still contribute to the diffusion
objective.  No LAM is executed online during NWM training.

We train Stage~1 for 60k steps with AdamW, learning rate
$2\times10^{-4}$, $(\beta_1,\beta_2)=(0.9,0.999)$, weight decay $0.01$, and
gradient clipping at 2.0.  Stage~2 replaces the latent adapter with a newly
initialized three-dimensional real-motion adapter for
$(\Delta x,\Delta y,\Delta\psi)$.
We train only this adapter for 3k steps and then jointly fine-tune it with the
world-model backbone for 100k steps; both use learning rate $10^{-4}$ and
AdamW with weight decay $0.01$.  Both stages use BF16 on eight NVIDIA L20 GPUs
with 48 examples per GPU, giving a global batch of 384 examples (1,536
context--target pairs per update).  We use the resulting 100k-step joint
checkpoint as the final OpenNWM model throughout; its internal experiment name
is \emph{OpenNWM-finalLAM-100k}.

For physical grounding, every target pose is expressed relative to the latest
context frame.  We divide planar translation by the dataset-specific waypoint
spacing---0.25 for RECON, 0.255 for HuRoN, 0.38 for SCAND, and 0.72 for
TartanDrive, all in meters per waypoint---and apply fixed min--max
normalization with
$x\in[-2.5,5]$, $y\in[-4,4]$, and
$\Delta\psi\in[-\pi,\pi]$.  Go Stanford uses 0.12 meters per waypoint at
evaluation time.  Thus, the implementation does not normalize translations by
an empirically estimated average step length.

\paragraph{Standalone Pixel-LAM implementation.}
The LAM reconstruction results use a Pixel-LAM trained independently of the
Pixel+Action LAM that supplies OpenNWM's offline proxies.  It processes
$240\times320$ RGB frames with $16\times16$ patches and represents a frame
transition with a 32-dimensional variational latent.  The architecture has a
24-block spatiotemporal encoder and a 24-block spatial decoder, both of width
1024 with 16 attention heads, for approximately 709M parameters.  Its loss is
\begin{equation}
    \mathcal{L}_{\mathrm{Pixel}}
    = \mathcal{L}_{\mathrm{RGB}}
    + 10^{-6}\mathcal{L}_{\mathrm{KL}},
\end{equation}
where $\mathcal{L}_{\mathrm{RGB}}$ is pixel-space mean squared error.

We train this model on 13 action-free sources together with RGB frames from
the RECON, HuRoN, and SCAND training splits; it does not use their
action annotations.  Training uses uniformly sampled signed offsets in
$[-8,8]$, AdamW with learning rate $2.5\times10^{-5}$ and weight decay 0.01,
a 12.8k-step linear warm-up followed by cosine decay to 0.1 of the peak rate,
and BF16 mixed precision.  We train for 100k optimizer steps on eight GPUs
with 20 examples per GPU (global batch 160).  This checkpoint is used only for
the public LAM reconstruction comparison and is distinct from the Pixel+Action
proxy used by OpenNWM.

\paragraph{Public LAM benchmark protocol.}
For each public baseline, we use its released checkpoint, native preprocessing,
and native RGB decoder.  We evaluate these baselines and our Pixel-LAM on
exactly the same record IDs, and compute all metrics through the same pipeline
on RGB in $[0,1]$ at $240\times320$ resolution.  The navigation-domain
aggregate contains the held-out RECON, HuRoN, and SCAND splits.  Go Stanford
is the transfer split and is unseen by our Pixel-LAM.
Because the full training exposure of every released checkpoint is not known,
we do not claim that these labels define universal ID and OOD splits for all
baselines.  For each dataset, we sample 100 valid pairs at each signed offset in
$\{-8,-6,-4,-2,2,4,6,8\}$, giving 2,400 navigation-domain pairs and 800 Go
Stanford pairs.
Offset zero is excluded.  Since several public LAMs do not release compatible
navigation-action decoders, this comparison is restricted to RGB
reconstruction.

\subsection{Comparison Results}

\begin{table*}[t]
  \centering
  \caption{\textbf{Main visual-prediction and navigation comparison.}
  Autoregressive prediction is evaluated on RECON, direct prediction on
  SACSoN, and navigation on RECON.  Baseline numbers are published references
  and are not all protocol-matched; see the text for the CEM settings.
  A dash marks a value unavailable in the audited results and does not support
  a ranking claim.  Lower is better for every reported metric.}
  \label{tab:main_prediction_navigation}
  \small
  \setlength{\tabcolsep}{6pt}
  \renewcommand{\arraystretch}{1.10}
  \begin{tabular}{@{}lcccccc@{}}
    \toprule
    \multicolumn{7}{@{}l}{\textbf{(a) Visual prediction}} \\
    \addlinespace[4pt]
    \multirow{2}{*}{Method}
      & \multicolumn{3}{c}{4\,s}
      & \multicolumn{3}{c}{16\,s} \\
    \cmidrule(lr){2-4}\cmidrule(lr){5-7}
      & LPIPS$\downarrow$ & DreamSim$\downarrow$ & FID$\downarrow$
      & LPIPS$\downarrow$ & DreamSim$\downarrow$ & FID$\downarrow$ \\
    \midrule

    \multicolumn{7}{@{}l}{\textit{Autoregressive on RECON}} \\
    \addlinespace[2pt]
    NWM~\cite{yang2026arforcing}
      & 0.423 & 0.181 & 63.0
      & 0.533 & 0.319 & 77.3 \\
    AR Forcing~\cite{yang2026arforcing}
      & 0.341 & 0.125 & 52.9
      & 0.463 & 0.210 & 66.0 \\
    \textbf{OpenNWM (Ours)}
      & -- & -- & --
      & -- & -- & -- \\

    \addlinespace[6pt]
    \multicolumn{7}{@{}l}{\textit{Direct on SACSoN}} \\
    \addlinespace[2pt]
    NWM~\cite{zhang2026raenwm}
      & 0.407 & 0.229 & 26.15
      & 0.470 & 0.281 & 33.06 \\
    RAE-NWM~\cite{zhang2026raenwm}
      & 0.303 & 0.145 & 15.09
      & 0.349 & 0.171 & 15.90 \\
    \textbf{OpenNWM (Ours)}
      & -- & -- & --
      & -- & -- & -- \\

    \midrule
    \multicolumn{7}{@{}l}{\textbf{(b) Navigation on RECON}} \\
    \addlinespace[4pt]
    Method
      & \multicolumn{3}{c}{ATE $\downarrow$}
      & \multicolumn{3}{c}{RPE $\downarrow$} \\
    \midrule
    GNM~\cite{bar2025navigation}
      & \multicolumn{3}{c}{$1.87{\pm}0.00$}
      & \multicolumn{3}{c}{$0.73{\pm}0.00$} \\
    NoMaD~\cite{bar2025navigation}
      & \multicolumn{3}{c}{$1.93{\pm}0.04$}
      & \multicolumn{3}{c}{$0.52{\pm}0.00$} \\
    NWM (published; CEM120)~\cite{bar2025navigation}
      & \multicolumn{3}{c}{$1.13{\pm}0.02$}
      & \multicolumn{3}{c}{$0.35{\pm}0.01$} \\
    LAM~\cite{garrido2026latentaction}
      & \multicolumn{3}{c}{1.40}
      & \multicolumn{3}{c}{0.40} \\
    \addlinespace[4pt]
    \textbf{OpenNWM (Ours; CEM80)}
      & \multicolumn{3}{c}{--}
      & \multicolumn{3}{c}{--} \\
    \bottomrule
  \end{tabular}
\end{table*}


\paragraph{Visual prediction.}
Table~\ref{tab:main_prediction_navigation}(a) summarizes autoregressive and
direct prediction at the reported horizons.  The former evaluates rollout
under feedback of generated observations, whereas the latter predicts a
distant target without recursively consuming model outputs.  Because the two
blocks use different datasets and the baseline rows are published references,
we interpret each block independently and do not infer that one prediction
mode is generally stronger than the other.  We make no quantitative ranking
claim because OpenNWM measurements under the table-specific protocols are not
yet available.

\paragraph{Navigation planning.}
Table~\ref{tab:main_prediction_navigation}(b) reports ATE and RPE.  For
OpenNWM, this experiment tests whether generated observations provide a
geometrically useful CEM planning signal.  The cited policy and world-model
rows retain their published settings, including the CEM120/CEM80 difference
noted above where applicable, and are contextual references unless an
explicitly matched result is provided.

\begin{table}[t]
  \centering
  \caption{\textbf{Standalone LAM RGB reconstruction.}
  The navigation-domain split aggregates held-out RECON, HuRoN, and SCAND
  (2,400 pairs); the Go Stanford transfer split contains 800 pairs and is
  unseen by our model.  Every method uses the same record IDs and nonzero
  signed offsets.  Bold and underline denote the best and second-best results.
  Our entry is the independently trained Pixel-LAM, not OpenNWM's Pixel+Action
  proxy.}
  \label{tab:lam_reconstruction}
  \small
  \setlength{\tabcolsep}{3.5pt}
  \renewcommand{\arraystretch}{1.08}
  \resizebox{\columnwidth}{!}{%
    \begin{tabular}{@{}lcccc@{}}
      \toprule
      \multirow{2}{*}{Method}
        & \multicolumn{2}{c}{Navigation domain}
        & \multicolumn{2}{c}{Go Stanford transfer} \\
      \cmidrule(lr){2-3}\cmidrule(l){4-5}
        & PSNR$\uparrow$ & LPIPS$\downarrow$
        & PSNR$\uparrow$ & LPIPS$\downarrow$ \\
      \midrule
      DiLA LAM~\cite{zhang2026dila}
                            & 13.180 & 0.5972 & 14.078 & 0.6211 \\
      Public DreamDojo LAM~\cite{gao2026dreamdojo}
                            & \underline{18.950} & 0.5292
                            & \underline{20.713} & \underline{0.4852} \\
      LARA/MoTo tokenizer~\cite{liu2026lara,chen2024moto}
                            & 15.752 & 0.5381 & 17.742 & 0.5267 \\
      Olaf-World LAM~\cite{jiang2026olaf}
                            & 14.764 & 0.6349 & 16.599 & 0.6142 \\
      villa-X LAM~\cite{chen2025villa0x0}
                            & 15.327 & \underline{0.5163} & 17.560 & 0.5499 \\
      \midrule
      \textbf{Ours (Pixel-LAM)}
                            & \textbf{21.071} & \textbf{0.5049}
                            & \textbf{21.912} & \textbf{0.4721} \\
      \bottomrule
    \end{tabular}%
  }
\end{table}


\paragraph{Visual reconstruction and generalization.}
Table~\ref{tab:lam_reconstruction} shows that our standalone Pixel-LAM obtains
the best PSNR and LPIPS among the evaluated models on both splits.  On the
navigation-domain aggregate, it achieves 21.071\,dB PSNR and 0.505 LPIPS: a
2.12\,dB gain over the public DreamDojo LAM and a 2.2\% relative LPIPS
reduction compared with the strongest public LPIPS baseline, villa-X.  On Go
Stanford, which is unseen by our model, it achieves 21.912\,dB and 0.472,
improving over the public DreamDojo LAM by 1.199\,dB and reducing LPIPS by
2.7\%.  The reconstruction advantage therefore persists on our held-out
transfer domain.

\subsection{Ablation Studies}

\begin{table*}[t]
  \centering
  \caption{\textbf{Available NWM pretraining-ablation results at 4\,s.}
  Single-step and rollouts are evaluated on RECON; OOD direct prediction is
  evaluated on Go Stanford.  The original-setup NWM row is a reference only;
  the three aligned rows share a downstream recipe.  Values are
  single-checkpoint point estimates.  The unfinished LatentPT and planning
  entries are omitted rather than represented by empty cells.  DS denotes
  DreamSim; bold denotes the best aligned result in each metric column.}
  \label{tab:pretraining_ablation_available}
  \footnotesize
  \setlength{\tabcolsep}{2.8pt}
  \renewcommand{\arraystretch}{1.15}
  \resizebox{\linewidth}{!}{%
    \begin{tabular}{@{}lcc*{12}{c}@{}}
      \toprule
      \multirow{3}{*}{Setup}
        & \multirow{3}{*}{Pretrain}
        & \multirow{3}{*}{Proxy}
        & \multicolumn{9}{c}{RECON}
        & \multicolumn{3}{c}{Go Stanford} \\
      \cmidrule(lr){4-12}\cmidrule(l){13-15}
        & &
        & \multicolumn{3}{c}{Single-step}
        & \multicolumn{3}{c}{Rollout (1\,FPS)}
        & \multicolumn{3}{c}{Rollout (4\,FPS)}
        & \multicolumn{3}{c}{Direct} \\
      \cmidrule(lr){4-6}\cmidrule(lr){7-9}\cmidrule(lr){10-12}\cmidrule(l){13-15}
        & &
        & LPIPS$\downarrow$ & DS$\downarrow$ & PSNR$\uparrow$
        & LPIPS$\downarrow$ & DS$\downarrow$ & PSNR$\uparrow$
        & LPIPS$\downarrow$ & DS$\downarrow$ & PSNR$\uparrow$
        & LPIPS$\downarrow$ & DS$\downarrow$ & PSNR$\uparrow$ \\
      \midrule
      NWM-Real (original setup) & No & N/A
        & 0.325 & 0.120 & 15.357
        & 0.385 & 0.152 & 14.579
        & 0.450 & 0.214 & 13.314
        & 0.653 & 0.462 & 11.202 \\
      \midrule
      Aligned & No  & N/A
        & 0.457 & 0.257 & 12.839
        & 0.454 & 0.254 & 13.198
        & \textbf{0.461} & 0.257 & \textbf{13.174}
        & 0.645 & 0.472 & 11.307 \\
      Aligned & Yes & None
        & \textbf{0.359} & \textbf{0.148} & \textbf{14.872}
        & \textbf{0.417} & \textbf{0.182} & 13.975
        & 0.488 & 0.258 & 12.555
        & \textbf{0.635} & \textbf{0.435} & \textbf{11.842} \\
      Aligned & Yes & Geometry
        & 0.362 & 0.153 & 14.841
        & \textbf{0.417} & 0.186 & \textbf{13.986}
        & 0.475 & \textbf{0.250} & 13.051
        & 0.639 & 0.439 & 11.587 \\
      \bottomrule
    \end{tabular}%
  }
\end{table*}


\paragraph{Pretraining and proxy conditioning.}
Table~\ref{tab:pretraining_ablation_available} reports only the currently
complete prediction rows under the earlier controlled protocol.  Relative to
the aligned model trained without pretraining, action-free time-only
pretraining improves all three metrics for single-step RECON prediction,
1-FPS RECON rollout, and Go Stanford direct prediction, but not for the 4-FPS
rollout.  Geometry conditioning does not consistently dominate time-only
conditioning.  The original-setup NWM row is included only as a reference and
is not a controlled comparison with the aligned rows.  Because the LatentPT
row is not yet populated, these results do not establish an advantage for
latent-action conditioning.

\paragraph{LAM training-recipe comparison.}
We study Pixel, Pixel+Action, Action-only, and DINO-feature reconstruction
recipes.  Pixel minimizes RGB reconstruction error, Action-only decodes a
three-dimensional relative-motion target, and DINO reconstructs frozen DINOv3
ViT-L/16 patch features.  Pixel+Action uses
\begin{equation}
    \mathcal{L}_{\mathrm{P+A}}
    = \mathcal{L}_{\mathrm{RGB}}
    + \lambda_{\mathrm{act}}\mathcal{L}_{\mathrm{act}}
    + 10^{-6}\mathcal{L}_{\mathrm{KL}},
\end{equation}
where the action term is masked on unlabeled examples.  Its mixed batches
contain 75\% action-free video and 25\% labeled navigation video, and
$\lambda_{\mathrm{act}}\approx0.172$ is calibrated once from the relative
initial magnitudes of the pixel and action losses.  The available checkpoints
do not form a strictly data-controlled objective ablation: Pixel and
Pixel+Action use the action-free sources together with the three physical
navigation datasets, whereas the audited Action-only and DINO checkpoints use
only the three physical datasets.  We therefore treat this as a comparison of
complete training recipes and do not attribute any difference solely to the
supervision objective.

\paragraph{Action-free data scale and label availability.}
Separately from the 100k-step Pixel-LAM benchmark and the 13-source OpenNWM
Stage-1 recipe, we conduct a controlled Pixel+Action data ablation on an
expanded 15-source action-free pool.  We vary the retained fraction of
trajectories as
$U\in\{25,50,75,100\}\%$ with all physical-action labels visible, and compare
$L\in\{0,50,100\}\%$ label availability at fixed $U=25\%$.  The $U$ subsets
are nested and split by trajectory after removing a shared held-out set.  The
$L$ setting changes only the action mask: RGB frames from RECON, HuRoN,
and SCAND continue to contribute to pixel reconstruction at every value of
$L$.

All six models use eight GPUs, 48 examples per GPU (global batch 384), 62.5k
optimizer steps, and therefore 24.0M sample presentations.  We use AdamW with
learning rate $4\times10^{-5}$, an 8k-step warm-up, and cosine decay to 0.1 of
the peak rate.  Checkpoints are evaluated every 10k steps and at 62.5k on
held-out trajectories from all 15 action-free sources, Go Stanford, and
TartanDrive.  Each dataset contributes 20 windows at offsets
$\{-8,-6,-4,-2,0,2,4,6,8\}$; the primary PSNR, SSIM, and LPIPS aggregates
exclude offset zero.  We defer empirical claims from this study until all six
runs have complete results.
